% explainer-source-hash: 1052a15d1085016c71a43388b5cf19496148b7f043cde42de9084a313d139b3f
% explainer-source-commit: cf2330e37a9828ebbd882ac081c6898966fee6b5
% explainer-generated: 2026-07-16
% Regenerate with /explain-all (see WIRING.md). Do not edit this stamp by hand:
% it is how the toolchain knows whether this document still matches the code.
\documentclass[11pt,a4paper]{article}
\input{preamble}

\title{\textbf{inference-engine}\\[0.3em]
\large An LLM serving engine built from scratch\\[0.6em]
\normalsize Brief}
\author{}
\date{}

\begin{document}
\maketitle
\thispagestyle{empty}

% ======================================================================
\section{The project in one page}

\file{inference-engine} turns a trained language model into a service. It is an
HTTP server that many clients can send prompts to simultaneously, that streams
words back as they are generated, and that shares one machine's memory between
all of them under a coherent policy.

It is Python and PyTorch, running GPT-2 family models (\code{gpt2},
\code{distilgpt2}) on CPU. Hugging Face supplies the trained weights and the
tokenizer. Everything that makes it a serving \emph{engine} is implemented in
the repository: the forward pass, the attention, the memory allocator, the
scheduler, the sampler, the speculative decoder, the HTTP surface.

\begin{keyidea}{Why a project like this exists}
Serving a language model is not a machine learning problem. It is an operating
systems problem wearing a machine learning hat: a memory allocator, a scheduler,
cache coherence, admission control, preemption. The three mechanisms that define
a modern serving engine (paged KV cache, continuous batching, speculative
decoding) are all systems techniques. This project rebuilds them small enough to
hold in one head, verifies them against a reference implementation, and measures
them honestly on the hardware actually available.
\end{keyidea}

\begin{table}[H]
\centering
\small
\begin{tabular}{lrl}
\toprule
\textbf{Area} & \textbf{Lines} & \textbf{Contents} \\
\midrule
\file{engine/} & 815 & cache, block manager, scheduler, model, sampling, speculative \\
\file{server/} & 206 & FastAPI, SSE, CLI \\
\file{tests/} & 731 & 7 modules, 55 tests \\
\file{benchmarks/} & 271 & harness plus committed JSON results \\
\bottomrule
\end{tabular}
\end{table}

% ======================================================================
\section{The three ideas, from first principles}

A language model can do exactly one thing: given some tokens, predict the next
one. Text is produced by looping: predict, append, predict again. Every word you
see from a chatbot cost one full pass through the model.

\subsection{Why a KV cache exists}

Each model layer performs \emph{attention}: every position looks back at every
earlier position, matching a \emph{query} against each one's \emph{key} to decide
how much of its \emph{value} to take. Predicting token 501 needs the keys and
values of tokens 1 through 500. Recomputing them each step makes generating $n$
tokens cost $O(n^2)$.

\begin{keyidea}{The cache, and the problem it creates}
A token's key and value depend only on the tokens before it. Attention is
\emph{causal}: the past cannot see the future, so the past never changes.
Compute each position's keys and values once, keep them, and every later step
reads them back. Generation becomes $O(n)$ and practical.

That store is the \textbf{KV cache}, and it converts repeated computation into
retained memory. For \code{gpt2} at this project's defaults the cache is about
604\,MB, against roughly 500\,MB of model weights. The weights are paid for once;
the cache is paid for per user, per token, forever. \textbf{The KV cache, not the
model, is what runs out first}, and managing it is what a serving engine is.
\end{keyidea}

\subsection{Idea one: page the cache}

The naive cache gives each sequence one contiguous region sized for the longest
answer it might produce. This over-reserves for tokens that may never exist, lets
one long request block short ones, and leaves freed gaps that are the wrong shape
for the next arrival.

\begin{figure}[H]
\centering
\resizebox{\textwidth}{!}{
\begin{tikzpicture}
  \node[lbl, font=\small\bfseries\upshape, text=warnred, anchor=west] at (0,1.0) {Naive: contiguous, sized for the worst case};
  \foreach \i/\name/\used in {0/A/3, 1/B/1, 2/C/2} {
    \node[lbl, anchor=east, font=\scriptsize] at (-0.15, -\i*0.8) {seq \name};
    \draw[draw=linegrey, fill=white] (0,{-\i*0.8-0.26}) rectangle (7,{-\i*0.8+0.26});
    \draw[draw=inkblue, fill=accent!30] (0,{-\i*0.8-0.26}) rectangle ({\used*0.85},{-\i*0.8+0.26});
    \node[font=\tiny, text=black!55] at ({\used*0.85+(7-\used*0.85)/2}, {-\i*0.8}) {reserved for tokens that may never exist};
  }
  \node[lbl, anchor=west, text=warnred, font=\scriptsize] at (0,-2.6) {over-reservation, head-of-line blocking, fragmentation};
  \begin{scope}[xshift=9cm]
  \node[lbl, font=\small\bfseries\upshape, text=okgreen, anchor=west] at (0,1.0) {Paged: fixed-size blocks, any order, on demand};
  \foreach \i/\name/\blocks in {0/A/{2,7,3}, 1/B/{5}, 2/C/{0,6}} {
    \node[lbl, anchor=east, font=\scriptsize] at (-0.15, {-\i*0.8}) {seq \name};
    \foreach \b [count=\j from 0] in \blocks {
      \draw[draw=inkblue, fill=accent!30] ({\j*0.95},{-\i*0.8-0.26}) rectangle ({\j*0.95+0.85},{-\i*0.8+0.26});
      \node[font=\tiny, text=inkblue] at ({\j*0.95+0.42},{-\i*0.8}) {blk \b};
    }
  }
  \node[lbl, anchor=west, text=okgreen, font=\scriptsize] at (0,-2.6) {waste $\le$ one partial block per sequence};
  \end{scope}
\end{tikzpicture}}
\caption{The reframing, borrowed from operating systems: a sequence stops owning
a \emph{region} and starts owning a \emph{list of block numbers}.}
\end{figure}

\begin{keyidea}{Paging, and the property that matters most}
Your operating system already does this. A program thinks it has one long stretch
of memory; the OS hands it scattered fixed-size \emph{pages} and keeps a
\emph{page table} translating what the program thinks it is using into where it
really is. Here: a sequence thinks it has smooth token positions, the engine
gives it scattered blocks (16 tokens each), and a \emph{block table} translates.

The immediate wins are that external fragmentation becomes impossible (every free
block is interchangeable) and waste is capped at one partial block per sequence.
But the \emph{valuable} consequence is subtler: once addresses are indirect, two
sequences can point at the same physical block. Sharing becomes a refcount bump.
Prefix caching and multi-completion forks both fall out of that for free.
\end{keyidea}

\begin{figure}[H]
\centering
\resizebox{0.85\textwidth}{!}{
\begin{tikzpicture}
  \node[lbl, anchor=west, font=\scriptsize\bfseries\upshape, text=inkblue] at (0,1.75) {logical positions the sequence believes in};
  \foreach \p in {0,...,7} {
    \draw[draw=linegrey, fill=softgrey] ({\p*0.8},1.0) rectangle ({\p*0.8+0.75},1.5);
    \node[font=\tiny, text=inkblue] at ({\p*0.8+0.37},1.25) {p=\p};
  }
  \node[lbl, anchor=west, font=\scriptsize\bfseries\upshape, text=inkblue] at (0,0.55) {\code{seq.block\_table}};
  \foreach \i/\b in {0/7, 1/2} {
    \draw[draw=inkblue, fill=accent!25, thick] ({\i*3.2},-0.35) rectangle ({\i*3.2+3.0},0.15);
    \node[font=\scriptsize, text=inkblue] at ({\i*3.2+1.5},-0.1) {table[\i] = block \b};
  }
  \node[lbl, anchor=west, font=\scriptsize\bfseries\upshape, text=inkblue] at (0,-0.9) {physical KV tensor, one flat row per slot};
  \foreach \b in {0,...,8} {
    \draw[draw=linegrey, fill=white] ({\b*0.8},-1.95) rectangle ({\b*0.8+0.75},-1.45);
    \node[font=\tiny, text=black!45] at ({\b*0.8+0.37},-1.7) {blk \b};
  }
  \foreach \b in {7,2} {
    \draw[draw=inkblue, fill=accent!25, thick] ({\b*0.8},-1.95) rectangle ({\b*0.8+0.75},-1.45);
    \node[font=\tiny, text=inkblue] at ({\b*0.8+0.37},-1.7) {blk \b};
  }
  \draw[flow] (1.5,-0.35) -- ({7*0.8+0.37},-1.45);
  \draw[flow] (4.7,-0.35) -- ({2*0.8+0.37},-1.45);
  \node[box, anchor=west, align=left, font=\scriptsize] at (0,-3.0)
    {\code{slot = block\_table[p // block\_size] * block\_size + p \% block\_size}\\
     with \code{block\_size = 4}: position 6 sits in \code{table[1] = block 2} at offset 2, flat slot 10.};
\end{tikzpicture}}
\caption{Address translation. One line of arithmetic, and everything else in the
memory system is a consequence of it.}
\end{figure}

Three features are built directly on the indirection:

\begin{itemize}
\item \textbf{Copy-on-write.} Two sequences share blocks until one writes, then
      the writer takes a private copy of just that block. Requesting $n$
      completions of one prompt costs one block copy each, not $n$ prompt
      recomputations. Full blocks stay shared forever: their contents can never
      change.
\item \textbf{Prefix caching.} Full blocks are content-hashed with a
      \emph{chained} SHA-256 (each block's hash includes the previous block's
      hash), so an identical token window under a different prefix cannot
      collide. A later request with the same prompt prefix re-references those
      blocks and skips computing those positions entirely. Freed hashed blocks
      go to an LRU evictable pool rather than the free list, so the cache costs
      nothing in capacity.
\item \textbf{Cheap preemption.} A preempted sequence loses its blocks and
      recomputes later, and the prefix cache usually hands most of them back.
\end{itemize}

\subsection{Idea two: batch continuously}

Running several sequences through the model together is much faster than one at a
time, because the weights are read from memory once and shared across the batch.
The naive approach freezes a batch until every member finishes; since they finish
at different times, most of the batch is computing padding for rows that are
already done, and a new request cannot start at all.

\begin{figure}[H]
\centering
\resizebox{\textwidth}{!}{
\begin{tikzpicture}[x=0.55cm, y=0.55cm]
  \node[lbl, font=\small\bfseries\upshape, text=warnred, anchor=west] at (0,1.4) {Static batching};
  \foreach \r/\name/\len in {0/A/12, 1/B/4, 2/C/6, 3/D/3} {
    \node[lbl, anchor=east, font=\scriptsize] at (-0.2,{-\r*0.8}) {\name};
    \draw[draw=inkblue, fill=accent!30] (0,{-\r*0.8-0.26}) rectangle (\len,{-\r*0.8+0.26});
    \ifnum\len<12
      \draw[draw=warnred, fill=warnred!12, dashed] (\len,{-\r*0.8-0.26}) rectangle (12,{-\r*0.8+0.26});
      \node[font=\tiny, text=warnred] at ({(\len+12)/2},{-\r*0.8}) {padding, computed and discarded};
    \fi
  }
  \draw[draw=warnred, very thick, dashed] (12,0.9) -- (12,-2.8);
  \node[lbl, anchor=west, text=warnred, font=\scriptsize] at (12.3,-1.2) {everyone waits for A;\\ a new request cannot join};

  \begin{scope}[yshift=-5.2cm]
  \node[lbl, font=\small\bfseries\upshape, text=okgreen, anchor=west] at (0,1.4) {Continuous batching: the batch is re-decided every token boundary};
  \foreach \r/\name/\e in {0/A/12, 1/B/4, 2/C/6, 3/D/3} {
    \node[lbl, anchor=east, font=\scriptsize] at (-0.2,{-\r*0.8}) {\name};
    \draw[draw=inkblue, fill=accent!30] (0,{-\r*0.8-0.26}) rectangle (\e,{-\r*0.8+0.26});
    \ifnum\e<12
      \node[lbl, anchor=west, text=okgreen, font=\tiny] at ({\e+0.2},{-\r*0.8}) {leaves, blocks freed immediately};
    \fi
  }
  \node[lbl, anchor=east, font=\scriptsize] at (-0.2,-3.2) {E};
  \draw[draw=okgreen, fill=okgreen!30] (2.4,{-3.2-0.26}) rectangle (10,{-3.2+0.26});
  \node[font=\tiny, text=inkblue] at (6.2,-3.2) {joins at the next token boundary};
  \foreach \t in {1,...,11} { \draw[draw=linegrey, dotted] (\t,1.1) -- (\t,-3.6); }
  \node[lbl, anchor=west, text=okgreen, font=\scriptsize] at (12.3,-1.6) {no padding, no idling,\\ blocks recycled continuously};
  \end{scope}
\end{tikzpicture}}
\caption{Where the project's measured 2.5x throughput gain comes from. Not from
making any request faster: from never computing padding and never idling. Dotted
lines are token boundaries, the only moments the scheduler may change its mind.}
\end{figure}

The scheduler (95 lines) admits FIFO while there is batch room and KV space, then
enforces one hard invariant: \textbf{the entire decode batch must be able to grow
by one token}.

\begin{keyidea}{Cumulative, not per-sequence: the bug worth knowing}
The obvious check tests each sequence against the free list independently. With
three sequences each needing one block and two blocks free, every individual
check passes, the batch is admitted, and the third sequence crashes mid-step
against an empty free list. The sequences are not competing with the free list;
they are competing with \emph{each other}, so the check must be summed across the
whole batch.

When the summed check fails, the newest-arrived sequence is preempted: its blocks
are freed and it is requeued \emph{at the front} of the waiting queue. Sacrificing
the newest protects the oldest, which has the most to lose; requeueing at the
front stops it being starved by later arrivals. Recovery is by recomputation
rather than swapping KV to slower memory, which on CPU would be a no-op with
extra bookkeeping.
\end{keyidea}

\subsection{Idea three: speculate}

\begin{keyidea}{Why guessing ahead can be free}
Decoding is \emph{memory-bound}. Producing one token means reading every weight
in the model out of memory, and the arithmetic units idle through most of it.
Running the model over five positions costs barely more than over one: the
weights are read once either way.

So let a small cheap model guess the next $k$ tokens, then run the big model
\emph{once} over all $k$ guesses and ask "would I have said that?" for each.
Every agreement is a token that cost a cheap forward instead of an expensive one.
\end{keyidea}

The subtlety is that this must not change what the user gets. The accept/reject
rule (Leviathan et al., 2023) guarantees it: accept a drafted token with
probability $\min(1, p/q)$ where $q$ is the draft's probability and $p$ the
target's; on the first rejection, sample a replacement from the residual
normalize($\max(p-q,0)$) and discard the rest. If the target liked the token less
than the draft did, keep it only $p/q$ of the time, which is exactly the fraction
that restores the target's rate; the residual is exactly the shape of the
resulting deficit. The two paths sum back to $p$ precisely.

\begin{figure}[H]
\centering
\resizebox{0.9\textwidth}{!}{
\begin{tikzpicture}[node distance=6mm]
  \node[box, align=left, fill=accent!8] (draft) {\textbf{Draft}: small model, \code{k} cheap forwards, one token each\\
    \scriptsize records token $x_i$ and full distribution $q_i$};
  \node[box, below=6mm of draft, align=left, fill=accent!8] (verify)
    {\textbf{Verify}: \textbf{ONE} big-model forward over all \code{k+1} positions $\rightarrow$ $p_0 \ldots p_k$};
  \node[dec, below=8mm of verify] (test) {$u < \min(1, p_i(x_i)/q_i(x_i))$?};
  \node[box, right=18mm of test, fill=okgreen!12, draw=okgreen] (acc) {\textbf{ACCEPT} $x_i$,\\ \scriptsize try $i+1$};
  \node[box, below=9mm of test, fill=warnred!10, draw=warnred, align=left] (rej)
    {\textbf{REJECT}: stop\\ \scriptsize replacement $\sim$ normalize($\max(p_i - q_i, 0)$)\\ \scriptsize discard later drafts, roll back BOTH caches};
  \node[box, left=18mm of test, fill=okgreen!12, draw=okgreen, align=left] (bonus)
    {\textbf{ALL ACCEPTED}\\ \scriptsize free bonus token from $p_k$\\ \scriptsize $k+1$ tokens, one big forward};
  \draw[flow] (draft) -- (verify);
  \draw[flow] (verify) -- (test);
  \draw[flow] (test) -- node[lbl, above] {yes} (acc);
  \draw[flow] (acc) to[bend right=40] node[lbl, right] {loop} (verify.east);
  \draw[flow] (test) -- node[lbl, right] {no} (rej);
  \draw[flow] (test) -- node[lbl, above] {loop ends} (bonus);
\end{tikzpicture}}
\caption{The accept/reject flow. The output is distributed \emph{identically} to
sampling from the target alone, not approximately. Speculative decoding trades no
quality for speed. At temperature 0 both distributions are one-hot and the rule
degenerates to "accept while the draft's argmax equals the target's".}
\end{figure}

% ======================================================================
\section{Architecture}

\begin{figure}[H]
\centering
\resizebox{0.95\textwidth}{!}{
\begin{tikzpicture}[node distance=6mm]
  \node[extbox, minimum width=9cm] (client) {HTTP clients (\code{curl}, \code{httpx}, OpenAI SDKs)};
  \node[box, below=7mm of client, minimum width=12.5cm, align=left] (app)
    {\textbf{server/app.py} \hfill FastAPI, OpenAI-compatible subset\\
     \scriptsize pydantic validation \quad SSE fan-out \quad one \code{asyncio.Queue} per sequence};
  \node[box, below=5mm of app, minimum width=12.5cm, align=left] (async)
    {\textbf{engine/async\_engine.py} \hfill AsyncEngine\\
     \scriptsize background step loop \quad blocking \code{step()} in a worker thread \quad one lock};
  \node[box, below=5mm of async, minimum width=12.5cm, align=left, minimum height=4.6cm, fill=white] (core) {};
  \node[anchor=north west, font=\small, text=inkblue] at ($(core.north west)+(4mm,-3mm)$)
    {\textbf{engine/llm\_engine.py} \hfill LLMEngine, the synchronous core: \code{step()}};
  \node[box, below left=11mm and -5.9cm of core.north] (sched)
    {\textbf{Scheduler}\\\scriptsize admission, batching\\\scriptsize preemption};
  \node[box, right=8mm of sched] (bm)
    {\textbf{BlockManager}\\\scriptsize free list, refcounts\\\scriptsize prefix hashes (LRU)\\\scriptsize returns COW ops};
  \node[box, right=8mm of bm] (spec)
    {\textbf{Speculative}\\\scriptsize draft model + own\\\scriptsize BlockManager, shadows};
  % 12mm, not 7mm: the edge label below needs room, or it lands on both boxes.
  \node[box, below=12mm of bm, minimum width=11cm] (runner)
    {\textbf{ModelRunner + GPT2} \file{engine/model.py}: paged attention over \file{cache.py} tensors};
  \draw[flow] (client) -- node[lbl, right] {POST /v1/completions (SSE)} (app);
  \draw[flow] (app) -- (async);
  \draw[flow] (async) -- (core.north);
  \draw[dflow] (sched) -- (bm);
  \draw[flow] (bm) -- node[lbl, right] {token ids, block tables, slots} (runner);
  \draw[flow] (spec) -- (runner);
\end{tikzpicture}}
\caption{Four layers, requests down and tokens up. \file{app.py} knows HTTP and
not blocks; \file{async\_engine.py} knows threads and not attention;
\file{llm\_engine.py} knows tokens and never touches a socket.}
\end{figure}

\begin{keyidea}{The best structural decision in the codebase}
\file{engine/block\_manager.py} does not import \code{torch}. It is pure
bookkeeping: it hands out integers and returns \emph{instructions} of the form
"copy physical block 4 to block 9". \file{engine/cache.py} is the only file that
touches KV tensors, and it obeys.

That is why the hardest logic in a serving engine is testable in isolation: 15
tests covering fragmentation, refcounting, copy-on-write, LRU eviction and
rollback run in well under a second because they never load a model. This is the
lesson from the project most worth reusing elsewhere.
\end{keyidea}

One \code{step()} advances every live sequence by one token boundary: schedule
(admitting arrivals, preempting under pressure), run each prefill alone, run all
decodes as one batched forward, then sample, detokenize and handle stopping.
Above it, \code{AsyncEngine} pushes the blocking \code{step()} onto a worker
thread via \code{asyncio.to\_thread} so the event loop stays free to stream, with
a single lock serializing engine mutations. The engine never learns what HTTP is;
it puts \code{StepOutput} objects into queues and \file{app.py} turns those into
SSE frames.

% ======================================================================
\section{How correctness is evidenced}

This is the strongest aspect of the project and its most transferable habit: the
claims are demonstrated at four different levels rather than asserted.

\begin{table}[H]
\centering
\small
\begin{tabular}{p{4.2cm}p{10.6cm}}
\toprule
\textbf{Level} & \textbf{Evidence} \\
\midrule
The model is right & \code{test\_parity}: this from-scratch GPT-2 forward versus Hugging Face's, largest disagreement across all 50{,}257 logits under \code{1e-3} (README reports about \code{3e-5} on distilgpt2, \code{2e-4} on gpt2: float32 noise). \\
The cache is right & \code{test\_paged\_forward\_matches\_cache\_free}: the paged path and the plain causal path agree to \code{1e-4}. \\
The allocator is right & \code{test\_scheduler\_invariants\_under\_churn}: 300 scheduling steps over six sequences in a deliberately undersized cache, asserting after \emph{every} step that \code{used + free == total}. A hand-rolled property test. \\
The speculation is right & \code{test\_empirical\_distribution\_matches\_target}: two tiny random transformers, 13-token vocabulary, 4000 full speculative rounds, asserting total variation distance from the target's true distribution under 0.05 (the noise floor at that sample count is about 0.02, and the test says so). Repeated with temperature and top-k applied. \\
\bottomrule
\end{tabular}
\end{table}

\begin{keyidea}{Verify against a reference before building on top}
Establishing logit parity first means every later bug is a \emph{systems} bug.
When the paged attention path misbehaved, the question was never "did I get the
transformer wrong?". That is why the parity test is the first one worth pointing
at, despite being the least interesting to read.

The speculative test is the sophisticated one. A probabilistic guarantee cannot
be checked by comparing two outputs, so the project shrinks the problem until
statistics become decisive: 13 tokens instead of 50{,}257, a two-layer random
transformer instead of GPT-2, 4000 samples. The threshold was derived from the
expected noise, not tuned until it passed.
\end{keyidea}

% ======================================================================
\section{Measurements}

All numbers are CPU-only on an Intel i7-1185G7 (8 threads), PyTorch float32,
greedy decoding. They are relative comparisons; they say nothing about absolute
serving performance. The committed JSON in \file{benchmarks/results/} records the
machine for each.

\subsection{What the project reports, and what a re-run found}

The committed numbers were re-measured on the same machine while this document was
written. That was possible only because the JSON records the host, an i7-1185G7,
and the model weights were cached locally. Conditions differed in three ways worth
stating: the re-run used \code{torch 2.10.0+cu128} and \code{transformers 5.14.1}
(the original's versions are not recorded), the host had a desktop session and
background processes, and PyTorch actually used \textbf{4} threads, since the
i7-1185G7 has 4 physical cores behind its 8 logical ones.

\begin{table}[H]
\centering
\small
\begin{tabular}{lrrl}
\toprule
\textbf{Claim the README makes} & \textbf{README} & \textbf{Re-run} & \textbf{Verdict} \\
\midrule
Throughput at concurrency 8 & 113.5 tok/s & 99.2 tok/s & 13\% lower \\
Continuous batching, 1 to 8 & 2.54x & \textbf{3.01x} & reproduced, stronger \\
Engine at 8 vs HF sequential & 2.84x & \textbf{3.11x} & reproduced, stronger \\
Speculative vs plain & 0.72x & \textbf{1.14x} & \textbf{did not reproduce} \\
Speculative acceptance rate & 0.5404 & \textbf{0.5404} & exact, to 4 decimals \\
\bottomrule
\end{tabular}
\end{table}

\begin{keyidea}{The headline claims hold, and the README is conservative}
\textbf{113.5 tok/s} did not reproduce exactly: this run measured 99.2 tok/s at
the same concurrency, on a contended host with unknown version drift. The figure
is the right order but should be quoted as "roughly 100 to 115 tok/s on a laptop
CPU depending on conditions", not as a constant.

\textbf{The 2.8x reproduced and came out at 3.11x.} The engine beat sequential
Hugging Face by more than claimed, because contention hurts a one-at-a-time
baseline more than a batched engine, which is precisely the effect continuous
batching exists to produce. The 2.5x batching claim likewise measured 3.01x. Both
headline claims are conservative on this machine, which is the right direction for
a claim to be wrong in.
\end{keyidea}

\begin{honest}{The speculative result reversed sign, and it is the real discrepancy}
The README's most distinctive claim is a negative one: speculative decoding is
"$\sim$0.7x slower here \ldots not a CPU speedup". The re-run measured
\textbf{1.14x faster}. Plain decode's wall time worsened 41\% (5.73s to 8.11s)
while speculative's \emph{improved} 11\% (7.96s to 7.11s).

This is not a correctness problem, and the evidence is unusually clean: 322 tokens
drafted and 174 accepted in both runs, to the token, six weeks and two library
versions apart. Greedy decoding is deterministic, so the algorithm made
bit-for-bit identical decisions; only the clock disagreed. The engine's
determinism claim is confirmed about as strongly as it could be.

A uniform slowdown preserves a ratio, and this one did not, so the effect is
differential. The coherent explanation is per-forward overhead: plain decode does
many tiny forwards (one token each), speculative does far fewer, fatter ones (k
draft forwards, then one target forward over k+1=5 positions). Anything raising
fixed per-forward cost, contention or a different torch build, penalizes the
many-small strategy and barely touches the fewer-large one. The same mechanism
explains why concurrency 1 fell to 0.74x of the original while concurrency 16 held
at 1.00x: big batches amortize that overhead, a batch of one cannot.

Ironically the README's own reasoning predicts this. It says speculative decoding
"pays off when the target forward dominates draft cost", then commits firmly to a
single sample of a quantity it had itself identified as conditional. The honest
position is narrower and stronger than either the README's or a naive "it is
actually faster": on this hardware at k=4 it sits near break-even, and which side
you land on depends on the software environment. Two measurements exist, 0.72x and
1.14x, and neither is a property of the algorithm. What is stable, and defensible
without qualification, is the acceptance rate and the exact output equivalence.
\end{honest}

\begin{honest}{The harness does not record the versions its numbers depend on}
\code{machine\_info()} captures CPU model, cores, platform, Python version and
\code{"device": "cpu"}. It records neither the \code{torch} version, the
\code{transformers} version, nor \code{torch.get\_num\_threads()}: precisely the
variables a throughput number depends on. The committed JSON therefore identifies
the \emph{machine} but does not permit reproducing the \emph{measurement}, which
is the point of committing it. It is why the discrepancy above cannot be fully
attributed: the experiment that would settle it is not recoverable from what was
recorded. Three more lines in \code{machine\_info()} would have prevented that
permanently.

(The committed JSON was restored byte-for-byte after the re-run; the repository's
recorded measurements are untouched.)
\end{honest}

% ======================================================================
\section{Honest assessment}

\subsection{What is genuinely strong}

The correctness evidence is real and layered rather than asserted. The
\code{BlockManager}-holds-no-tensors split makes the hardest logic testable in
under a second. Every figure in the README transcribes faithfully from the
committed JSON, with no rounding in the flattering direction. And the project
reports its own headline feature going backwards on CPU, with the reason, rather
than quietly dropping it: a README that says 0.72x is more persuasive than one
claiming 2x.

\begin{keyidea}{The pattern behind every serious finding below}
All of them live in the space \emph{around} the algorithms rather than in them.
The paged cache, the scheduler, the accept/reject sampler: correct, well tested,
defensible in detail. The failures are in dependency declaration, process
lifetime, error propagation and cross-component invariants, exactly the set of
things unit tests do not reach and a short benchmark cannot reveal.

That is a useful thing to be able to say about one's own work: the hard part was
done well and the gaps are in the boring parts. It is also the honest answer to
"what next", ahead of the GPU kernel.
\end{keyidea}

\subsection{The findings that matter}

The top three were found by \emph{running} the project, not reading it. The code
reads well enough that none of them are visible from the source alone.

\begin{honest}{Any exception in \code{step()} permanently bricks the server (verified)}
\code{AsyncEngine.\_run\_loop} has no exception handling at all. It runs as a
bare \code{asyncio.Task}; nothing awaits it, attaches a done-callback, or
inspects its exception. So when \code{step()} raises, the loop dies silently and
the server goes on accepting requests it will never answer, while
\code{/healthz} keeps returning \code{ok}.

Driven end-to-end through \code{AsyncEngine}, the stream hung for 25 seconds with
no tokens, the task was \code{done()} with
\code{RuntimeError('draft KV cache exhausted')}, and then a subsequent, entirely
innocent \code{"Hello there"} request \emph{also} hung: one bad request bricks
the process until someone restarts it. This ranks first because it is a severity
multiplier, converting every other bug (including ones not yet found) from a bad
response into a total outage. The fix is routine: wrap the step, fail the
affected sequences, keep looping.
\end{honest}

\begin{honest}{The draft cache's capacity invariant is false and reachable (verified)}
\file{docs/ARCHITECTURE.md} states that the draft cache's "block usage is always a
subset of the target's; one capacity check then covers both caches", and a
matching comment sits in \file{llm\_engine.py}. The scheduler indeed checks only
the target's block manager.

The reasoning covers \emph{retention} (the draft returns freed blocks immediately
rather than parking them in an evictable pool), and that much is true. It omits
\emph{sharing}, which runs the other way: the target's \code{fork()} path lets
sequences share physical blocks, while the draft allocates private blocks for
every shadow. The draft therefore needs \emph{more} blocks than the target, the
opposite of the claimed subset relation. Note this does not even depend on prefix
caching, so the stated \emph{reason} ("runs without prefix caching") does not
deliver the invariant on its own terms.

Driven with \code{num\_blocks=12}, a 63-token shared prompt and \code{n=4}: step
one showed four sequences sharing four blocks (the sharing working exactly as
designed), the target-side check passed comfortably, and step two died with
\code{RuntimeError: draft KV cache exhausted} because the draft needed sixteen
blocks against twelve. Combined with the finding above, that is a remotely
triggerable permanent outage.

At the default configuration arithmetic saves it by luck, not by the invariant:
\code{max\_batch\_size} 8 times \code{max\_model\_len} 1024 is exactly the 8192
token slots the cache holds. Raise the batch size and the margin is gone; the
project's own benchmark harness starts its server with \code{--max-batch-size 16}
(though never with a draft model). This is the most serious \emph{design} finding
because it is a documented invariant that the code does not have: answering "why
is one check enough?" from the architecture document would be giving a wrong
answer.
\end{honest}

\begin{honest}{\file{requirements.txt} declares a range the project cannot run in (verified)}
\file{engine/model.py} calls \code{from\_pretrained(model\_name, dtype=...)}. The
\code{dtype} keyword is recent; Hugging Face spelled it \code{torch\_dtype} for
years. \file{requirements.txt} declares \code{transformers>=4.44,<6}. Bisected
against the real package index: 4.51.3 and 4.55.4 both fail with
\code{TypeError: GPT2LMHeadModel.\_\_init\_\_() got an unexpected keyword
argument 'dtype'} and load no model at all, while 4.56.0 and 5.14.1 work. The
true floor is \textbf{4.56.0}.

So every version from 4.44 to 4.55.4, roughly two years of releases and the
larger part of the declared range, yields a clone that cannot load a single
model. This surfaced on the machine used for this document, which had 4.51.3
installed system-wide: the entire model-backed suite errored out. The fix is
\code{transformers>=4.56,<6}. The finding is not the severity but the category:
the project's central claim is that nothing is asserted that was not verified by
running it, and its declared environment is not the environment it was verified
in.
\end{honest}

\begin{honest}{A verified memory leak in the long-running part}
\code{LLMEngine.add\_request} inserts a \code{\_SeqState} per sequence.
\code{abort()} removes it for sequences still waiting; \code{\_finish()}, the
path every normally completing sequence takes, never does. So a server retains
one entry per request served, forever, each holding that request's complete
generated text. Confirmed by driving it, not by reading: after twenty sequential
generations the scheduler's queues were empty and the block manager was back to a
full free list, while \code{seq\_states} still held twenty entries. Nothing
catches it because every test and benchmark constructs an engine, runs a few
requests, and exits: precisely the blind spot of a project about long-running
services.
\end{honest}

\begin{honest}{Seeded determinism has a load-dependent hole for \code{n > 1}}
When a sibling completion cannot join the batch, the code has \emph{already}
consumed a random draw from its generator sampling a token it then discards
before requeueing. The same request with the same seed therefore produces
different output depending on whether the batch happened to be full. Seeded
determinism is an advertised feature; the tests use a batch size that never
triggers the branch, and the multi-completion HTTP test checks that two choices
come back but never that they reproduce.
\end{honest}

\begin{honest}{The attention is not paged at the kernel level}
The README discloses this and the code confirms it exactly: \code{gather} pulls
each sequence's scattered KV into a \emph{dense padded} tensor and calls a
standard attention kernel. The \textbf{memory} savings are entirely real; the
\textbf{compute} savings are not, since a batch mixing a 500-token sequence with
seven 20-token ones pads all eight to 500. The gather also materializes a copy of
the whole batch's context every layer, every step. A real engine writes a fused
kernel that walks the block table inside attention. The stated reason for not
doing so (on CPU with a small model it was not the bottleneck) is defensible, and
this is the biggest gap between the project and the systems it studies, so it is
worth volunteering rather than being caught by.
\end{honest}

\begin{honest}{The TTFT curve is a readout of unbatched prefill}
Time to first token goes 0.069s, 0.266s, 0.541s, 1.266s at concurrency 1, 4, 8,
16: near-linear. Prefills run one at a time, and a single prefill of these prompts
takes about 0.07s, so eight simultaneous arrivals means eight sequential prefills
and the last client waits for all of them ($8 \times 0.07 \approx 0.55$s against a
measured 0.541s). The README calls the growth "the expected tradeoff" and
separately lists unbatched prefill under future work, without connecting the two.
The curve is not a property of continuous batching; batching ragged prompts would
flatten most of it. \emph{(This arithmetic is inferred from the committed JSON and
the code, not separately measured.)}
\end{honest}

\begin{honest}{Smaller findings, in descending order}
The incremental detokenizer decodes the \emph{entire} output every step and
slices off what was already sent, making it $O(n^2)$ in output length (invisible
at the benchmark's 32 tokens, half a million token-decodes at GPT-2's 1024 limit).
\code{\_match\_prefix} recomputes its whole SHA-256 chain three times per
admission. \code{EOS\_TOKEN\_ID = 50256} is hardcoded at module level rather than
read from the tokenizer that the engine already holds, making the engine
silently GPT-2-family-only in a way the configuration does not express.
\code{num\_speculative\_tokens} is applied in \code{add\_request} even with no
draft model configured, while \code{Scheduler.lookahead} correctly is not: two
places disagreeing about the same question. \code{k} is a minimum across the
batch, so one nearly-finished sequence throttles everyone's lookahead. A
just-admitted prompt can be preempted in the same \code{schedule()} call that
admitted it (correct, but wasted work). And the prefix cache trusts SHA-256
without ever comparing tokens on a hit: the right call, but its correctness rests
on a cryptographic assumption rather than a check.
\end{honest}

\subsection{What the benchmarks do and do not support}

\begin{honest}{"2.8x over baseline" is honest, conservative, and not controlled}
It compares 113.5 tok/s (this engine, 8 concurrent, driven over HTTP with SSE
parsing) against 40.0 tok/s (HF \code{generate()}, in-process, no serialization).
Both are real and both are in the committed JSON, but they come from different
harnesses. The bias favours the baseline, so the true gap is if anything wider,
which is the right way to be wrong. The honest framing is "2.8x end-to-end
including our serving overhead against a bare in-process loop". Note too that
sequential HF \code{generate()} is a weak baseline by construction; the strong
one, HF with its own batching, is not measured.

Two further caveats: the sweep's harness counts \emph{SSE events}, not tokens
(one event is one token for plain decoding, so the totals are right, but the same
harness run with a draft model would understate throughput by up to 5x with no
warning); and the speculative benchmark measures a warm prefix cache, since it
warms up on the same four prompts it then times. The latter is fair between the
two engines it compares, which is what that benchmark is for.
\end{honest}

\begin{honest}{The gap the project cannot close, and says so}
There is no comparison against a production engine like vLLM, because vLLM needs
a GPU and this machine has none. The README states this as absent, with the
reason, rather than fabricating a number: the right call, since a made-up figure
would be worse than a missing one. But it means the project can demonstrate that
its \emph{mechanisms} work and cannot demonstrate that they work \emph{well}.
"Continuous batching gives 2.5x on 8 CPU threads" is what was measured; nothing
here says how the design behaves where these designs actually live.
\end{honest}

\subsection{The design choices worth defending as made}

\begin{itemize}
\item \textbf{One coarse engine lock.} Correct here: the forward dominates step
      time so completely on CPU that contention is unmeasurable, and one lock is
      the version a person can reason about. It stops being right the moment you
      want to overlap prefill with decode or pipeline the sampler.
\item \textbf{Recompute-only preemption.} On CPU the "device" and host are the
      same memory, so swapping KV out would be a no-op with extra bookkeeping.
\item \textbf{Unbatched prefill.} Defensible on CPU (prefill is already
      compute-bound), and the honest cost is the TTFT curve above.
\item \textbf{Speculative decoding kept despite going backwards.} It is
      correctness-verified and 0.72x on CPU, because the draft's forward is not
      cheap relative to the target's. It pays off when the target dominates draft
      cost (large model, GPU). Measuring it honestly and finding it negative is
      more useful than assuming the paper's GPU numbers transfer, and that is the
      right thing to say about it.
\end{itemize}

\end{document}
